\section{Transport recovery and the order of limits}\label{sec:inverse-limits}

\subsection{The observed invariant}
In this inverse problem the compact family and its unitary
one-parameter bundle class are known. The normalization
$L=\log(1/|t|)$ and the three component areas $A_0,A_1,A_2$ are known.
The parameter $\theta$, the transition matrices and their cycle
product are not supplied. The observations are the five lowest
eigenvalues of the actual operator $\mathcal L_{t,\theta}$, including
its zero modes. The target is the conjugacy class of
$\Ad(U_\theta)$ in $\operatorname{SO}(3)$.
It is a global bundle-transport invariant. The integral Chern number
computed in Section~\ref{sec:geometry} is not an output of this
spectral inverse.

\begin{theorem}\label{thm:transport-inverse}
With the above observations define
\begin{equation}\label{eq:transport-estimator}
 \widehat s_t=\frac{A_0A_1^3A_2}{8}
     \left(\frac{L}{2\pi}\right)^5\prod_{k=0}^4\lambda_k(t,\theta).
\end{equation}
For the error in Theorem~\ref{thm:five-spectrum},
\begin{equation}\label{eq:inverse-relative}
 (1-\varepsilon(t))^5\sin^2\theta
 \le\widehat s_t\le(1+\varepsilon(t))^5\sin^2\theta.
\end{equation}
Thus $3-4\widehat s_t$ recovers $\tr\Ad(U_\theta)$ with error
$O(\varepsilon(t)\sin^2\theta)$. If
$\operatorname{clip}(x)=\max(0,\min(1,x))$, then
\begin{equation}\label{eq:angle-estimator}
 \widehat\alpha_t=2\arcsin\sqrt{\operatorname{clip}(\widehat s_t)}
\end{equation}
converges uniformly over the family to its principal projective
conjugacy angle in $[0,\pi]$.
On a class $\sin^2\theta\in[\kappa,1-\kappa]$, $0<\kappa<1/2$,
the angle error is at most $C_{A,\kappa}(\varepsilon(t)+\delta)$
if each observed rescaled eigenvalue $L\lambda_k/(2\pi)$ has
additive error at most $\delta$, for small $\delta$ and $t$.
\end{theorem}

\begin{proof}
In coordinates $(a,x,y,z,b)$ write $c_\theta=\cos\theta$ and
$s_\theta=\sin\theta$. The energy and mass matrices are
\begin{equation}\label{eq:five-matrices}
 Q_\theta=\begin{pmatrix}
 2&-2&0&0&0\\
 -2&4&0&0&-2c_\theta\\
 0&0&4&0&-2s_\theta\\
 0&0&0&4&0\\
 0&-2c_\theta&-2s_\theta&0&2
 \end{pmatrix},\qquad
 M_A=\operatorname{diag}(A_0,2A_1,2A_1,2A_1,A_2).
\end{equation}
The Schur complement of the middle $4I_3$ block is
$\left(\begin{smallmatrix}1&-c_\theta\\-c_\theta&1\end{smallmatrix}\right)$.
Consequently
\begin{equation}\label{eq:determinant-product}
 \det Q_\theta=64\sin^2\theta,\qquad
 \prod_{k=0}^4\mu_k(\theta)
       =\frac{8\sin^2\theta}{A_0A_1^3A_2}.
\end{equation}
Multiplying the five relative bounds in
\eqref{eq:five-relative} proves \eqref{eq:inverse-relative}.
At $\theta\in\pi\ZZ$ both products are exactly zero.

Since $U_\theta=\operatorname{diag}(e^{i\theta},e^{-i\theta})$,
its adjoint rotation has trace
$1+2\cos(2\theta)=3-4\sin^2\theta$. The principal angle is
$2\arcsin\sqrt{\sin^2\theta}$. Uniform continuity on $[0,1]$
proves uniform recovery after clipping. On the stated interior class
the derivative of this function is bounded, which proves Lipschitz
dependence on the estimator.

The rescaled eigenvalues lie in a fixed bounded interval by
\eqref{eq:five-relative}. Their product is a polynomial with bounded
derivative on a fixed bounded enlargement. Errors at most $\delta$
in the five factors therefore change \eqref{eq:transport-estimator}
by at most $C_A\delta$. Combine this with the preceding angle bound.
\end{proof}

No uniform Lipschitz angle bound is asserted at the endpoint angles.
In particular, fixed absolute noise need not permit relative recovery
of an arbitrarily small angle. The observations do not distinguish a
signed $\SU(2)$ lift, central scalar twists, or reversal of the cycle
orientation. Nor does this one-parameter statement imply spectral
rigidity among arbitrary bundles outside the specified family.

\begin{corollary}\label{cor:weakest-rate}
Let $\vartheta$ be the small signed distance from $\theta$ to
$\pi\ZZ$. Jointly as $t\to0$ and $\vartheta\to0$ with
$\vartheta\ne0$,
\begin{equation}\label{eq:weakest-rate}
 \lambda_0(t,\theta)=
 \frac{2\pi\vartheta^2}{L(A_0+2A_1+A_2)}(1+o(1)).
\end{equation}
At $\vartheta=0$ the eigenvalue is exactly zero.
\end{corollary}
\begin{proof}
At $\theta=0$ the kernel is spanned by
$w=(1,1,0,0,1)^{\mathsf T}$. Direct cofactors give
$\operatorname{adj}(Q_0)=64ww^{\mathsf T}$.
The coefficient of $-\mu$ in
$\det(Q_0-\mu M_A)/\det M_A$, equal to the product of its four
positive generalized eigenvalues, is therefore
\[
 \frac{\tr(\operatorname{adj}(Q_0)M_A)}{\det M_A}
       =\frac{8(A_0+2A_1+A_2)}{A_0A_1^3A_2}.
\]
The simple small eigenvalue is analytic and even in $\theta$:
reflection of the $y$ coordinate identifies the matrices at
$\theta$ and $-\theta$. Divide \eqref{eq:determinant-product}
by the analytic nonzero product of the other four eigenvalues to get
\[
 \mu_0(\theta)=\frac{\theta^2}{A_0+2A_1+A_2}+O(\theta^4).
\]
A sign change of the $b$ coordinate reduces the calculation at any
other multiple of $\pi$ to this one. The relative, parameter-uniform
bound in Theorem~\ref{thm:five-spectrum} proves the joint assertion.
\end{proof}

The symbolic determinant and cofactor identities in this proof are
reproduced by the accompanying computation. They check the finite
matrix normalization; they do not replace the proof relating its
eigenvalues to the actual threefold operator.

\subsection{A common space for the time limits}
To compare metrics on different fibres, pull them to the cover and use
the fixed component frames on compact subsets of $S\times D_j$
away from the nodes. All limits below mean $L^2$ convergence on every
such compact subset, relative to $h_S\otimes I$.
An effective tuple $P$ denotes the corresponding relative metrics
$I_{T^{1,0}S}\otimes P_j$. The reconstruction differs from its
constant tuple by $O(L^{-1})$ on these subsets. Global $L^2$
comparison implies this component convergence.
Write $H_\theta(\tau;P)$ for the effective flow in rescaled time
$\tau=s/L$, initialized at $P$. Its equation is independent of $t$.

\begin{proposition}\label{prop:time-limits}
For matching initial data
$h_{t,\theta}(0)=\widehat K_{t,\theta}(P)$ the following hold.
\begin{enumerate}
\item For fixed $\theta\notin\pi\ZZ$ and an invariant tuple
$P\ne I$, physical time and degeneration do not commute:
\begin{equation}\label{eq:physical-noncommutation}
 \lim_{s\to\infty}\lim_{t\to0}h_{t,\theta}(s)=P,
 \qquad
 \lim_{t\to0}\lim_{s\to\infty}h_{t,\theta}(s)=I.
\end{equation}
\item On a fixed rate class $\mu_0(\theta)\ge\kappa>0$, for
fixed $\theta$ and $P$, the limit at each $\tau\ge0$ is
$H_\theta(\tau;P)$. The limits $t\to0$ and $\tau\to\infty$
commute and equal $I$. More generally, for bounded tuples and every
sequence $t\to0$, $\tau(t)\to\infty$, with $\theta$ in this
rate class, $h_{t,\theta}(L\tau(t))\to I$.
\item If $\theta(t)\to0$ with $\theta(t)\notin\pi\ZZ$, take
$P_j=\exp(a\sigma_x)$ for all $j$, with $a\ne0$. Then
\begin{equation}\label{eq:weak-rate-noncommutation}
 \lim_{\tau\to\infty}\lim_{t\to0}
       h_{t,\theta(t)}(L\tau)=P,
 \qquad
 \lim_{t\to0}\lim_{\tau\to\infty}
       h_{t,\theta(t)}(L\tau)=I.
\end{equation}
\end{enumerate}
\end{proposition}

\begin{proof}
For a fixed physical time $s$, effective speed $O(L^{-1})$ makes the
effective tuple tend to $P$. In \eqref{eq:bounded-tracking} take
$\eta=0$ and $T=s/L$. Its constant is independent of $T$, so the
error is at most $Cs(1+\log L)^2/L^2$, which tends to zero.
This proves the first limit in \eqref{eq:physical-noncommutation}.
On each fixed stable smooth fibre, the relative logarithm has no
parallel trace-free kernel. The coercivity argument in the proof of
Theorem~\ref{thm:all-time}, using the fixed-fibre positive gap,
implies convergence to $h_0$. This proves the other limit.

On the rate class in part 2, the all-time comparison has error tending
to zero uniformly, and the effective trajectory tends to $I$
exponentially in $\tau$ with uniform constants on bounded tuples.
These facts prove both the fixed-$\tau$ limit and the uniform diagonal
statement. They also prove the claimed interchange of the two limits.

At $\theta=0$ the tuple in part 3 is stationary: $R=I$ and all
its entries coincide. It is invariant under $J$ and is a global
parallel metric at that parameter. On each bounded $\tau$ interval,
smooth dependence of the confined effective vector fields on $\theta$
gives $H_{\theta(t)}(\tau;P)\to P$ uniformly. The bounded-time
actual-flow estimate is uniform in $\theta$, hence the actual
metrics have the same limit. Each fixed nonzero $t$, on the other
hand, has a stable parameter and converges to $I$ at infinite time.
This proves \eqref{eq:weak-rate-noncommutation}.
\end{proof}

\begin{corollary}\label{cor:spatial-limits}
For the entire bounded open class of Theorem~\ref{thm:spatial-entry},
uniform diagonal relaxation to $I$ holds when
$\mu_0(\theta_t)\ge\kappa>0$, $t\to0$ and $\tau(t)\to\infty$.
No convergence assumption on the selected tuples $P_t$ is needed for
this conclusion. If in addition $P_t\to P$ and $\theta_t\to\theta$
within this rate class, then for each fixed $\tau>0$ the component
limit is $H_\theta(\tau;P)$. Its infinite-$\tau$ limit is $I$,
and the corresponding iterated limits commute.
\end{corollary}
\begin{proof}
The selected tuples lie in a common bounded set, the all-time error
in \eqref{eq:spatial-alltime} tends to zero, and $s_*/L\to0$.
Uniform effective-flow decay gives the diagonal assertion. Under the added
phase and parameter convergence assumptions, finite-dimensional continuous
dependence identifies the limit at each positive $\tau$; entry has
then occurred for all small $t$. The remaining assertion follows
from effective-flow decay and fixed-fibre convergence.
\end{proof}

A general spatial initial metric at $\tau=0$ need not equal its
selected phase. Nor need a nonconvergent family of selected phases
have any fixed-$\tau$ degeneration limit. Finally, part 3 of
Proposition~\ref{prop:time-limits} shows failure of uniform relaxation
to $I$ near reducibility. It does not disprove tracking by a
parameter-dependent effective trajectory, whose own relaxation time
also diverges there.
